\documentclass[11pt]{report}

\usepackage[a4paper,margin=1in]{geometry}
\usepackage{amsmath,amssymb,amsthm,mathtools,bm}
\usepackage{booktabs}
\usepackage{enumitem}
\usepackage{parskip}
\usepackage{microtype}
\usepackage[dvipsnames,table]{xcolor}
\usepackage{tikz}
\usepackage{pgfplots}
\pgfplotsset{compat=1.18}
\usetikzlibrary{arrows.meta,positioning,calc}
\usepackage[breakable]{tcolorbox}
\usepackage[colorlinks=true,linkcolor=NavyBlue,citecolor=NavyBlue,urlcolor=NavyBlue]{hyperref}

% ---------- math shortcuts ----------
\DeclareMathOperator*{\argmax}{arg\,max}
\DeclareMathOperator*{\argmin}{arg\,min}
\newcommand{\E}{\mathbb{E}}
\newcommand{\R}{\mathbb{R}}
\newcommand{\Teams}{\mathcal{T}}
\newcommand{\Probs}{\mathcal{P}}
\newcommand{\Conts}{\mathcal{C}}
\newcommand{\Obs}{\mathcal{O}}
\newcommand{\given}{\,\vert\,}

% ---------- theorem-like ----------
\theoremstyle{definition}
\newtheorem{definition}{Definition}[chapter]
\newtheorem{remark}{Remark}[chapter]
\newtheorem{exercise}{Exercise}[chapter]
\numberwithin{equation}{chapter}

% ---------- callout boxes ----------
\newtcolorbox{keyidea}{breakable,colback=NavyBlue!6,colframe=NavyBlue!55!black,
  fonttitle=\bfseries,title=Key idea,left=2mm,right=2mm,top=1mm,bottom=1mm}
\newtcolorbox{intuition}{breakable,colback=ForestGreen!7,colframe=ForestGreen!55!black,
  fonttitle=\bfseries,title=Intuition,left=2mm,right=2mm,top=1mm,bottom=1mm}
\newtcolorbox{worked}[1]{breakable,colback=BurntOrange!7,colframe=BurntOrange!65!black,
  fonttitle=\bfseries,title=Worked example: #1,left=2mm,right=2mm,top=1mm,bottom=1mm}
\newtcolorbox{takeaway}{breakable,colback=black!4,colframe=black!45,
  fonttitle=\bfseries,title=Chapter in one paragraph,left=2mm,right=2mm,top=1mm,bottom=1mm}

\title{\textbf{Rating Problems from Standings}\\[4pt]
\large A Gentle, Practical Guide for ICPC-Style Contests}
\author{Bilguun}
\date{\today}

\begin{document}
\maketitle

\tableofcontents

\chapter*{Before we begin}
\addcontentsline{toc}{chapter}{Before we begin}

You have the standings of many programming contests. For each contest you know
which teams competed, which problems each team solved, and at what minute each
solve happened. You want one number per problem: \emph{how hard was it?} On
Codeforces this number comes for free, because every contestant already has a
rating. In ICPC there is no such rating. This guide builds the missing
machinery from the ground up.

The plan is deliberately gentle. We start from a single, very intuitive idea ---
\emph{a problem's difficulty is the skill level at which solving it is a
coin flip} --- and we keep returning to it. Every formula is introduced first in
words, then in symbols, and then exercised on a tiny five-team contest that we
carry through the whole book so the numbers stay concrete. By the end you will
have two complete methods, a way to use the solve-time data that most tools throw
away, and a principled recipe for putting problems from different regions of the
world onto one comparable scale.

No background beyond first-year probability and a little calculus is assumed.
Familiarity with the Elo rating system or with logistic regression will make some
passages feel like old friends, but it is not required.

\chapter{The problem, and the one idea behind everything}

\section{Why the obvious measures fail}

Suppose problem~D in some regional was solved by only $3$ of $200$ teams. Is it
hard? Probably --- but the solve count alone cannot tell you, because it ignores
\emph{who} solved it. Three solves by the three weakest teams would be
astonishing and would mean the problem is easy; three solves by the three
strongest teams is exactly what you would expect from a hard problem. The
acceptance rate has the same blind spot. Any honest difficulty measure has to
weigh each solve by the strength of the team that produced it.

That single sentence already tells us we cannot rate problems without first
having some notion of team strength --- and, awkwardly, we cannot measure team
strength without knowing how hard the problems they solved were. The two
quantities define each other. Resolving that circularity is the technical heart
of this guide; everything else is bookkeeping around it.

\begin{keyidea}
The difficulty of a problem is defined relative to \emph{solvers}, not to
\emph{counts}. A problem is exactly as hard as the skill level at which a
contestant has a $50\%$ chance of solving it. Stronger solvers shift that
level up; weaker solvers shift it down.
\end{keyidea}

\section{The coin-flip definition}

Fix a scale on which both ``team ability'' and ``problem difficulty'' are just
real numbers, measured in the same units (we will borrow the Codeforces scale, so
the numbers feel familiar). Write $\theta$ for a team's ability and $b$ for a
problem's difficulty. The promise of the scale is this:

\begin{center}
\emph{a team whose ability equals a problem's difficulty has a $50\%$ chance of
solving it.}
\end{center}

A team that is much stronger than the problem ($\theta \gg b$) almost always
solves it; a team much weaker ($\theta \ll b$) almost never does. Between those
extremes the probability slides smoothly. The smooth curve that competitive
programming has settled on is the logistic curve of the next chapter, but the
\emph{definition} of difficulty does not depend on the exact curve: $b$ is
wherever the probability passes through one half.

\begin{intuition}
Think of difficulty as a height of a high-jump bar and ability as how high a
jumper can clear. Set the bar at the height the athlete clears half the time and
you have measured the athlete; raise the bar until even the best clears it only
half the time and you have measured the bar. Problems and teams are measured on
the same ruler precisely because they trade places in the same probability.
\end{intuition}

\section{Our running example}

We will use one toy contest for the rest of the book. Five teams, whose
(pretend, known) abilities on the Codeforces scale are listed below, attempt four
problems labelled A--D. A check means the team solved the problem during the
contest.

\begin{table}[h]
\centering
\renewcommand{\arraystretch}{1.2}
\begin{tabular}{lc cccc}
\toprule
Team & Ability $\theta$ & A & B & C & D\\
\midrule
$T_1$ & $2400$ & \checkmark & \checkmark & \checkmark & \\
$T_2$ & $2200$ & \checkmark & \checkmark & & \\
$T_3$ & $2000$ & \checkmark & \checkmark & & \\
$T_4$ & $1800$ & \checkmark & & & \\
$T_5$ & $1600$ & \checkmark & & & \\
\midrule
\multicolumn{2}{r}{solved by} & $5$ & $3$ & $1$ & $0$\\
\bottomrule
\end{tabular}
\caption{The running example. Problem A was solved by everyone, D by no one,
and B and C land in between. In a real project the abilities would be
\emph{unknown} and estimated; here we fix them so we can check the arithmetic by
hand.}
\label{tab:toy}
\end{table}

Intuitively A is very easy, D very hard, B moderately hard (only the top three
got it), and C hard (only the very top team got it). Our job over the next
chapters is to turn that intuition into numbers, and then --- the real trick ---
to recover the abilities themselves when they are \emph{not} handed to us.

\begin{takeaway}
We cannot rank problems by how many teams solved them, because that ignores how
strong those teams were. Instead we define a problem's difficulty as the ability
level at which a team is equally likely to solve it or not. Ability and
difficulty live on one shared ruler and are defined in terms of each other, which
is the circularity the rest of the book resolves.
\end{takeaway}

\chapter{The logistic building block}

\section{From ``coin flip'' to a curve}

We need an actual formula for the probability that a team of ability $\theta$
solves a problem of difficulty $b$. Three things are non-negotiable: the
probability is $\tfrac12$ when $\theta=b$, it rises toward $1$ as the team
outclasses the problem, and it falls toward $0$ in the opposite case. The
logistic (sigmoid) function does exactly this, and it is the same function the
Elo system uses for chess. On the Codeforces scale,
\begin{equation}\label{eq:pi}
  \pi(\theta,b)\;=\;\frac{1}{1+10^{\,(b-\theta)/400}}.
\end{equation}
The constant $400$ sets how quickly the curve turns: a $200$-point edge gives
about a $76\%$ solve chance and a $400$-point edge about $91\%$. Figure~\ref{fig:sigmoid}
plots it against the \emph{gap} $\theta-b$.

\begin{figure}[h]
\centering
\begin{tikzpicture}
\begin{axis}[
  width=11cm, height=6cm,
  domain=-600:600, samples=120,
  xlabel={ability minus difficulty, $\theta-b$},
  ylabel={solve probability $\pi$},
  ymin=0, ymax=1, xmin=-600, xmax=600,
  xtick={-400,-200,0,200,400},
  ytick={0,0.25,0.5,0.75,1},
  grid=both, grid style={black!10},
  axis lines=left,
  tick label style={font=\small},
  label style={font=\small},
]
\addplot[NavyBlue,very thick] {1/(1+10^(-x/400))};
\addplot[only marks,mark=*,mark size=1.6pt,BurntOrange]
  coordinates {(0,0.5) (200,0.76) (400,0.91)};
\draw[black!35,dashed] (axis cs:0,0) -- (axis cs:0,0.5);
\draw[black!35,dashed] (axis cs:-600,0.5) -- (axis cs:0,0.5);
\end{axis}
\end{tikzpicture}
\caption{The logistic solve curve~\eqref{eq:pi}. At an equal match the
probability is exactly one half (the dashed cross-hair); the orange dots mark the
$+200$ and $+400$ edges.}
\label{fig:sigmoid}
\end{figure}

It is sometimes cleaner to write \eqref{eq:pi} with the natural logistic
$\sigma(x)=1/(1+e^{-x})$,
\begin{equation}\label{eq:sigmascale}
  \pi(\theta,b)=\sigma\!\Big(\frac{\theta-b}{s}\Big),
  \qquad s=\frac{400}{\ln 10}\approx 173.7,
\end{equation}
which makes plain that only the \emph{gap} $\theta-b$ matters, rescaled by $s$.
Keep both forms in mind: \eqref{eq:pi} for hand arithmetic on the Codeforces
scale, \eqref{eq:sigmascale} for the modelling chapters.

\begin{intuition}
Every $400$ points of advantage multiply a team's \emph{odds} of solving (the
ratio $\pi/(1-\pi)$) by ten. That is the entire content of the constant in the
exponent: rating differences are a logarithmic currency, and $400$ points buys
one order of magnitude.
\end{intuition}

\section{The same curve rates head-to-head matches}

The logistic also gives the probability that one team outperforms another: team
$u$ beats team $t$ with probability $\pi(\theta_u,\theta_t)$, reading the
opponent's ability in place of a problem's difficulty. We will need this in
Chapter~4 to score how a team actually \emph{placed} in a contest. The two uses
--- ``did you solve the problem?'' and ``did you beat the rival?'' --- are the
same equation wearing different labels.

\begin{takeaway}
The probability of solving is a logistic function of the gap between ability and
difficulty: one half at parity, climbing one order of magnitude in odds per $400$
points. The identical curve, with an opponent's rating in place of a problem's,
scores head-to-head outcomes.
\end{takeaway}

\chapter{Rating a problem when the teams are known}

This chapter does the easy half of the circular problem: it assumes the team
abilities are already known and rates the problems. (Chapter~4 removes the
assumption.) This is exactly what tools such as \texttt{clist.by} compute for
Codeforces, where the abilities are supplied by the judge.

\section{Expected solvers versus actual solvers}

Pick a candidate difficulty $b$. For each team we can compute its solve
probability $\pi(\theta_t,b)$, and summing these gives the number of solvers we
would \emph{expect} at that difficulty:
\begin{equation}\label{eq:expected}
  \E[\text{solvers}](b)\;=\;\sum_{t}\pi(\theta_t,b).
\end{equation}
Now turn it around. We do not know $b$, but we \emph{do} know the actual number
of solvers $S$. So we choose the difficulty at which the expected count matches
the observed count.

\begin{keyidea}
A problem's difficulty is the value $b$ that makes the expected number of solvers
equal the number who actually solved it:
\[
  \sum_t \pi(\theta_t,b)\;=\;S.
\]
Because the left side decreases steadily as $b$ grows (harder problems are solved
by fewer), there is exactly one such $b$, found by binary search.
\end{keyidea}

\begin{definition}[Weighted rating]\label{def:wr}
Given teams with abilities $\{\theta_t\}$ and nonnegative weights $\{w_t\}$ (take
all $w_t=1$ for now; Section~\ref{sec:weight} explains the weights), and a target
$S$, define $\mathrm{WR}(\{(w_t,\theta_t)\},S)$ as the unique $b$ solving
$\sum_t w_t\,\pi(\theta_t,b)=S$.
\end{definition}

\begin{worked}{difficulty of problem B}
Problem B was solved by $3$ teams. Using the abilities from
Table~\ref{tab:toy}, we look for $b$ with $\sum_t\pi(\theta_t,b)=3$. Trying
$b=2000$ gives expected solvers
\[
  0.91+0.76+0.50+0.24+0.09 = 2.50,
\]
which is below $3$, so the true difficulty is a little \emph{lower}. Trying
$b=1900$ gives $0.95+0.85+0.64+0.36+0.15=2.95$, just under $3$. A few more
bisection steps land near $b\approx 1890$. So B's difficulty is about
$1890$ --- between $T_3$ ($2000$) and $T_4$ ($1800$), which matches the fact that
exactly the top three teams solved it.
\end{worked}

\begin{worked}{difficulty of problem C}
Problem C was solved by $1$ team. We seek $b$ with $\sum_t\pi(\theta_t,b)=1$.
At $b=2200$ the expected count is about $1.62$ (too many), at $b=2400$ about
$0.87$ (too few), and around $b\approx 2360$ it equals $1$. So C rates near
$2360$, just below the lone solver $T_1$ at $2400$ --- a genuinely hard problem.
\end{worked}

\section{The two ends of the scale}

Problem A was solved by \emph{everyone} and D by \emph{no one}. The equation
$\sum_t\pi=S$ has no finite solution in these cases: to expect all five solvers
the difficulty must fall to $-\infty$, and to expect zero it must rise to
$+\infty$. In practice you floor and cap the scale (Codeforces effectively treats
``solved by all'' as rating $0$ and ``solved by none'' as a large ceiling). The
lesson is not a defect but a caution: a problem at the very edge of a field is
only weakly measured --- all you learn is ``easier than the weakest team'' or
``harder than the strongest team here.'' To pin it down you need a field that
straddles it.

\section{Weighting unreliable teams}\label{sec:weight}

Not every ability estimate deserves equal trust. A team you have seen in only one
contest could be a fluke; a team with a long history is a steady yardstick. The
weights $w_t$ in Definition~\ref{def:wr} let steadier teams speak louder. A simple,
battle-tested choice (used by \texttt{clist}) grows the weight from near zero
toward one as a team accumulates a history of $n_t$ prior contests:
\begin{equation}\label{eq:weight}
  w_t = 1 - 0.9^{\,n_t+1}.
\end{equation}
A newcomer ($n_t=0$) gets weight $0.1$; after ten contests the weight is about
$0.65$; a veteran approaches $1$. The same idea will reappear as \emph{priors} in
Chapter~5, in a more principled dress.

\begin{takeaway}
If you already know team abilities, a problem's difficulty is whatever value
makes the summed solve probabilities equal the actual solver count --- a quantity
found by binary search. Problems solved by all or none sit at the floor or
ceiling and are only loosely measured, and shaky teams should be down-weighted so
their noise does not leak into the difficulty.
\end{takeaway}

\chapter{Breaking the circle: rating teams and problems together}

Chapter~3 assumed the abilities. In ICPC we do not have them, so we must estimate
abilities and difficulties at the same time. The escape from the circularity is
an old trick: guess one side, use it to compute the other, and repeat until the
two stop changing.

\section{Scoring how a team placed}

Before we can update an ability we need to score a team's contest result. The
natural score is the rating consistent with where it \emph{finished}. If a team
placed $r$-th, then $r-1$ teams beat it, so the ability we infer is the one whose
\emph{expected} number of superiors equals $r-1$. Using the head-to-head form of
the logistic, the team's \emph{performance rating} $\rho$ in that contest solves
\begin{equation}\label{eq:perf}
  1+\sum_{u\neq t}\pi(\theta_u,\rho)\;=\;r .
\end{equation}
This is the very same weighted-rating inversion as before, now with rank in place
of solver count.

\begin{worked}{a team that overperformed}
Suppose in our toy field team $T_4$ (true ability $1800$) finished \emph{2nd}.
Then the expected number of teams above it must equal $1$, i.e.\ we solve
$\sum_{u\neq T_4}\pi(\theta_u,\rho)=1$ over the other four teams
($2400,2200,2000,1600$). Searching, $\rho\approx 2350$: finishing second in this
field is performance worthy of a $2350$ player, far above $T_4$'s usual $1800$.
A good night, and the system records it as such.
\end{worked}

\section{The alternating loop}

Now the engine. Start every team at a prior guess. Sweep through all contests,
computing each team's performance from the \emph{current} abilities of its rivals.
Then move each team's ability toward the average of its performances across all
the contests it played. Repeat. When nobody's rating moves appreciably, stop and
rate the problems with the converged abilities.

\begin{figure}[h]
\centering
\begin{tikzpicture}[
  box/.style={rounded corners,draw=NavyBlue!60!black,fill=NavyBlue!6,
    align=center,inner sep=6pt,minimum width=4.2cm,minimum height=1.5cm,font=\small},
  >={Stealth[length=3mm]}]
\node[box] (A) at (0,0) {\textbf{Team abilities} $\theta_t$\\(current guess)};
\node[box] (B) at (8,0) {\textbf{Performances} $\rho_{t,c}$\\(one per contest played)};
\draw[->,thick] (A.north) to[bend left=22]
  node[above,font=\small\itshape]{score each result} (B.north);
\draw[->,thick] (B.south) to[bend left=22]
  node[below,font=\small\itshape]{average back into ability} (A.south);
\end{tikzpicture}
\caption{The alternating fixed point. Abilities produce performances; performances
are averaged back into abilities, and the loop repeats until the abilities stop
changing --- at which point the problems are rated with the converged abilities.
The fixed point is a set of ratings internally consistent with every contest at once.}
\label{fig:loop}
\end{figure}

The update that closes the loop blends a team's performances with its prior, with
the experience weights of \eqref{eq:weight}:
\begin{equation}\label{eq:update}
  \theta_t \;\leftarrow\;
  \frac{\displaystyle\sum_{c\,:\,t\text{ played }c} w_{t,c}\,\tfrac12\big(\rho_{t,c}+\theta_t^{\text{prior}}\big)}
       {\displaystyle\sum_{c\,:\,t\text{ played }c} w_{t,c}} .
\end{equation}

\begin{intuition}
Why does this resolve the circularity? Because each pass uses the \emph{best
current} estimate of the other side and pushes both estimates toward mutual
agreement. It is the same logic as repeatedly averaging to find an equilibrium:
the abilities that explain the standings, given difficulties that explain the
standings, given those abilities. A team that played several contests is the glue
--- its single ability has to fit all of them, which drags those contests onto a
common footing. Hold onto that observation; in Chapter~7 it becomes the entire
basis for comparing different regions.
\end{intuition}

This method (essentially what \texttt{clist} does, plus the ability layer it
normally gets for free) is quick to build and easy to sanity-check. Its limits
are that it returns bare numbers with no error bars, and a problem with very few
solvers can swing wildly on the luck of a single strong-but-underrated team.
Chapter~5 fixes both.

\begin{takeaway}
When abilities are unknown, alternate: score each team by the rating its finishing
place implies, average those performances back into its ability, and iterate to a
fixed point; then rate the problems. Teams that appear in several contests force
those contests into a single consistent scale.
\end{takeaway}

\chapter{A cleaner foundation: item response theory}

The alternating loop works, but it is a procedure rather than a model: it does
not say what we believe about the data, only what steps to run. Stating an honest
probabilistic model buys us uncertainty estimates, graceful handling of the
all/none edge cases, and a clean way to fold in priors. The right model is one
that education researchers have used for decades to grade test questions ---
\emph{item response theory}.

\section{The Rasch model}

Treat each ``did team $t$ solve problem $p$?'' as a coin flip whose bias is set by
the gap between ability and difficulty. That is the Rasch model, and it is just
our logistic curve read as a likelihood:
\begin{equation}\label{eq:rasch}
  \Pr(\text{team }t\text{ solves }p)\;=\;\pi(\theta_t,b_p)
  =\sigma\!\Big(\frac{\theta_t-b_p}{s}\Big).
\end{equation}
A student is to a test what a team is to a contest: able people pass hard items
only sometimes, weak people pass easy items most of the time, and every response
is one logistic coin flip. Estimating all the $\theta$'s and $b$'s at once is the
same statistical task whether the ``items'' are exam questions or contest
problems.

\begin{intuition}
The alternating loop of Chapter~4 was secretly trying to fit this model; it just
went back and forth instead of optimizing everything together. Writing the model
down lets us hand the whole thing to a single optimizer and ask not only for the
best estimates but for how sure we are of them.
\end{intuition}

\section{The likelihood, and why shared teams link contests}

Assuming the solves are independent given the abilities and difficulties, the
probability of the entire observed dataset is the product over all observed
(team, problem) pairs $\Obs$, and its logarithm is
\begin{equation}\label{eq:loglik}
  \ell(\bm\theta,\bm b)=\sum_{(t,p)\in\Obs}
  \Big[y_{tp}\log\pi(\theta_t,b_p)+(1-y_{tp})\log\big(1-\pi(\theta_t,b_p)\big)\Big],
\end{equation}
where $y_{tp}\in\{0,1\}$ records the solve. The crucial structural fact: the
\emph{same} symbol $\theta_t$ appears in every term for every contest team $t$
entered. A team that played a regional and the World Finals ties those two
contests' problems together automatically, with no separate rescaling step,
because one ability must explain both. This is the model-level version of the
``glue'' from Chapter~4, and it is what Chapter~7 exploits.

\section{Discrimination, priors, and uncertainty}

Some problems split the field sharply --- strong teams solve them, weak teams do
not, with little middle ground --- while others are noisy for everyone. A
two-parameter extension gives each problem a \emph{discrimination} $a_p>0$
controlling how steep its curve is:
\begin{equation}\label{eq:twopl}
  \Pr(\text{solve})=\sigma\!\big(a_p(\theta_t-b_p)\big).
\end{equation}
A large $a_p$ is a sharp filter; a small $a_p$ a near-coin-flip regardless of
skill.

Small regionals have few teams, so some problems and teams are barely observed.
Left alone, the likelihood can send a never-solved problem's difficulty to
infinity. The cure is a \emph{prior}: a gentle prior belief that abilities and
difficulties are not absurdly extreme,
\begin{equation}\label{eq:priors}
  \theta_t\sim\mathcal N(\mu_t,\sigma_\theta^2),\qquad
  b_p\sim\mathcal N(\mu_b,\sigma_b^2).
\end{equation}
We then maximize the log-likelihood plus the log-prior (the \emph{maximum a
posteriori}, or MAP, estimate):
\begin{equation}\label{eq:map}
  (\hat{\bm\theta},\hat{\bm b})=\argmax_{\bm\theta,\bm b}\;
  \ell(\bm\theta,\bm b)
  -\frac{1}{2\sigma_\theta^2}\sum_t(\theta_t-\mu_t)^2
  -\frac{1}{2\sigma_b^2}\sum_p(b_p-\mu_b)^2 .
\end{equation}
The prior is exactly the experience weight of Chapter~4 grown up: instead of
shrinking a shaky team's \emph{influence}, it shrinks a shaky estimate toward a
sensible default, and it does so for problems too. A problem nobody solved now
contributes a term that pushes its difficulty up but is held finite by the prior,
so the infinity of Section~3.2 never arises. And because we have a full model,
fitting it in a Bayesian way (by MCMC or variational inference) returns not just a
number per problem but a credible interval --- precisely what you want before
trusting a difficulty estimated from four solvers.

\begin{takeaway}
Reading the logistic curve as a likelihood gives the Rasch model from education
testing: every solve is one coin flip set by ability minus difficulty. Fitting all
parameters jointly links any contests that share a team, a discrimination
parameter lets problems differ in sharpness, and Gaussian priors both tame
tiny-field noise and yield honest uncertainty.
\end{takeaway}

\chapter{Getting more from the clock}

Standings carry a signal most difficulty tools discard: \emph{when} each solve
happened. A problem solved at minute~15 was easy for that team; the same problem
solved at minute~280, in the dying moments, was a struggle. We can fold this in
without leaving the logistic family.

\section{Solving as a race against the clock}

Model each team's attack on each problem as a race: at every instant the team has
some rate of cracking the problem, higher when the team outclasses the problem.
Concretely, let the instantaneous solve rate (the \emph{hazard}) be
\begin{equation}\label{eq:hazard}
  \lambda_{tp}=\lambda_0\,\exp\!\Big(\frac{\theta_t-b_p}{s}\Big).
\end{equation}
A strong team facing an easy problem has a high rate and tends to solve quickly; a
weak team facing a hard problem has a low rate and usually runs out of time. Two
kinds of observation feed the fit. A team that solved $p$ at time $\tau_{tp}$
contributes the probability density of solving \emph{exactly then}; a team that
never solved $p$ contributes the probability of \emph{not yet} having solved it
when the clock $T_c$ ran out --- a \emph{censored} observation, in the language of
survival analysis. Their log-likelihood terms are
\begin{equation}\label{eq:surv}
  \ell_{tp}=
  \begin{cases}
    \log\lambda_{tp}-\lambda_{tp}\,\tau_{tp}, & \text{solved at }\tau_{tp},\\[4pt]
    -\,\lambda_{tp}\,T_c, & \text{never solved.}
  \end{cases}
\end{equation}

\begin{intuition}
The binary model of Chapter~5 only asked \emph{whether} a team solved a problem.
The clock model also asks \emph{how fast}, which is extra evidence about
difficulty for free. And the ``never solved'' case is handled honestly: it does
not mean the team \emph{couldn't}, only that it didn't within the time limit --
exactly what a censored race says.
\end{intuition}

A pleasant bonus: contest length $T_c$ enters the formula explicitly, which
partly corrects the unfairness of comparing a five-hour regional against a
two-hour short round --- the same problem looks easier when teams had more time,
and the $T_c$ term knows it.

\begin{takeaway}
Treat each solve as a timed race whose speed grows with the ability-minus-difficulty
gap. Solved problems contribute when they were solved; unsolved ones contribute as
races that hadn't finished at the buzzer. This uses the solve-time data, handles
``nobody solved it'' gracefully, and bakes in a correction for contest length.
\end{takeaway}

\chapter{Putting different regions on one scale}

So far every method produces ratings that are internally consistent. But
``internally consistent'' has a hidden loophole that matters enormously once we
compare regions.

\section{The floating-zero problem}

Look again at the model: the solve probability depends only on the
\emph{difference} $\theta_t-b_p$. So if we add $100$ to every team's ability and
$100$ to every problem's difficulty, every probability --- and therefore the
entire likelihood --- is completely unchanged. The data cannot distinguish the
two solutions. Within a single contest this is harmless: it just means ``the zero
of the scale is arbitrary.'' Across two contests with no common teams and no
common problems, it is fatal: each contest fixes its own arbitrary zero, and a
$2000$ in one may be a $2300$ in the other. Their difficulties are simply not
comparable.

\begin{keyidea}
Ratings are only pinned down up to a shared additive shift. Two contests are
comparable only if something forces their shifts to agree --- and the only thing
that can is a participant (or, rarely, a problem) they have in common.
\end{keyidea}

\section{The linking graph}

This turns standardization into a question about connectivity. Build a graph whose
nodes are contests; join two contests by an edge when they share enough teams.
Two contests can be placed on the same scale exactly when a path of shared-team
edges connects them. A shared team carries one ability into both contests'
likelihoods, nailing down their relative shift; chains of shared teams propagate
the alignment outward.

\begin{figure}[h]
\centering
\begin{tikzpicture}[
  reg/.style={rounded corners,draw=ForestGreen!60!black,fill=ForestGreen!8,
    align=center,inner sep=5pt,minimum width=2.7cm,minimum height=1.1cm,font=\small},
  hub/.style={rounded corners,draw=NavyBlue!65!black,fill=NavyBlue!10,
    align=center,inner sep=5pt,minimum width=2.7cm,minimum height=1.1cm,font=\small\bfseries},
  >={Stealth[length=2.6mm]},
  every edge/.style={draw=black!45,thick}]
\node[reg] (r1) at (0,2.1) {Asia\\regional};
\node[reg] (r2) at (0,0) {Europe\\regional};
\node[reg] (r3) at (0,-2.1) {Latin\\regional};
\node[hub] (wf) at (4.2,1.1) {World Finals};
\node[hub] (uc) at (4.2,-1.1) {Universal Cup};
\node[reg,fill=black!4,draw=black!45] (scale) at (8.4,0) {one shared\\scale};
\foreach \r in {r1,r2,r3}{
  \draw (\r) edge[->] (wf);
  \draw (\r) edge[->] (uc);
}
\draw[->] (wf) -- (scale);
\draw[->] (uc) -- (scale);
\node[font=\scriptsize\itshape,text=black!55] at (2.1,2.15) {shared teams};
\end{tikzpicture}
\caption{Each regional is an island whose internal ratings float on their own
zero. The World Finals (qualifying teams from every region meet on one problem
set) and the Universal Cup (recurring shared rounds across regions) are the
bridges whose shared teams fuse the islands into a single scale.}
\label{fig:link}
\end{figure}

Two structures do the heavy lifting. The \textbf{World Finals} gathers qualifying
teams from every region onto one problem set, so each finalist carries its
regional ability into a shared arena. The \textbf{Universal Cup} runs the same
rounds for many strong teams across regions, repeatedly --- a far denser web of
shared participations than the once-a-year Finals, and in practice the stronger
linker of the two. Before trusting any cross-region number, compute the connected
components of this graph: anything joined to the Finals/Universal-Cup component is
comparable; anything cut off from it must be reported on its own scale, with a
clear warning.

\section{Anchoring the zero, and a free cold start}

Connectivity removes all the floating zeros but one: the single global shift for
the whole connected pool. Fix it with one convention. The plain choice is to
declare the average ability to be zero. The better choice ties the scale to
numbers people already understand --- the Codeforces ratings of the team members.
For a team whose members carry Codeforces ratings, combine those member ratings
into a team prior and use it both to seed the optimization and to anchor the
scale:
\begin{equation}\label{eq:cfprior}
  \mu_t=\mathrm{WR}\big(\{(1,R_m)\}_{m\in t},\,\tfrac12\big),
\end{equation}
where $R_m$ are the members' Codeforces ratings and $\mathrm{WR}$ is the same
inversion from Definition~\ref{def:wr}. This one move does three jobs at once: it
gives brand-new teams a sensible starting ability (no cold-start problem), it
fixes the global shift, and it makes every output readable as a
Codeforces-equivalent number --- so ``this Yokohama problem is about $2400$''
means what a competitive programmer expects it to mean.

\begin{takeaway}
Because only ability-minus-difficulty is identified, each isolated contest has its
own arbitrary zero, and cross-region comparison requires a chain of shared teams
to align those zeros. The World Finals and especially the Universal Cup supply
that chain; one global convention --- ideally anchoring to members' Codeforces
ratings --- fixes the last remaining shift and gives the whole scale familiar
units.
\end{takeaway}

\chapter{A recipe, and the traps along the way}

\section{Build it in two stages}

Start simple and earn the complexity. \textbf{Stage one:} run the alternating
loop of Chapter~4, seeded and anchored by the Codeforces bootstrap of
\eqref{eq:cfprior}. It is a day's work on top of the open-source \texttt{clist}
code and gives ratings you can eyeball immediately. Validate them against
problems whose difficulty you already have an opinion on --- World Finals problems
with known editorial ratings, or problems mirrored on a rated judge. \textbf{Stage
two:} once you trust the linking graph's connectivity, move the final numbers to
the Bayesian item-response model of Chapter~5, with the solve-time likelihood of
Chapter~6 if you want the extra signal and the error bars. The first stage tells
you whether the data and the plumbing are sound; the second gives you the ratings
you publish.

\section{Three traps}

\begin{remark}[A team is not a stable thing]
Rosters change between seasons, so ``University X, 2024'' and ``University X,
2025'' are usually different people and should not share one ability. Where you
know the rosters, model individual contestants and combine them into a team
strength with \eqref{eq:cfprior}; this is both more accurate and a bonus for
linking, since one person recurs across years and shows up in the Universal Cup
with different teammates, adding edges to the graph of Chapter~7.
\end{remark}

\begin{remark}[Spelling is load-bearing]
The linking graph is only as good as your identifiers. If a team is written one
way in a regional's standings and another way at the Finals or on a mirror, the
shared-team edge silently vanishes and your careful standardization quietly
breaks. Reconciling names and identities across sources is not clean-up work to do
later; it is a precondition for everything in Chapter~7.
\end{remark}

\begin{remark}[Few solvers, and unequal clocks]
A problem solved by a handful of teams is estimated from a handful of coin flips;
treat its rating as a wide interval, not a point, which is one more reason to
prefer the Bayesian fit for the final numbers. And remember that a five-hour
regional and a two-hour round are not the same yardstick --- the solve-time model
of Chapter~6 helps, but cross-format comparisons always deserve a second look.
\end{remark}

\section*{Exercises}
\addcontentsline{toc}{section}{Exercises}

\begin{exercise}
Using the abilities in Table~\ref{tab:toy}, verify by bisection that problem~B's
difficulty is near $1890$. Then recompute it after deleting the weakest team
$T_5$ from the field. Does the difficulty move, and why?
\end{exercise}

\begin{exercise}
A problem is solved by exactly the top half of a field of evenly spaced abilities.
Argue, without heavy computation, that its difficulty should land near the
\emph{median} ability, and connect this to the coin-flip definition of
Chapter~1.
\end{exercise}

\begin{exercise}
Show that adding a constant $\kappa$ to every $\theta_t$ and every $b_p$ leaves
the likelihood~\eqref{eq:loglik} unchanged. Explain in one sentence why this makes
an isolated contest's zero arbitrary, and what removes the arbitrariness across
contests.
\end{exercise}

\begin{exercise}
In the clock model~\eqref{eq:surv}, suppose you doubled every contest's time limit
$T_c$ but the underlying hazards stayed the same. Qualitatively, which way would
estimated difficulties move, and is that the behaviour you want?
\end{exercise}

\chapter*{Where these ideas come from}
\addcontentsline{toc}{chapter}{Where these ideas come from}

The logistic solve curve and the $0.5/200/400$ calibration are the Codeforces
rating and problem-difficulty conventions. The expected-versus-actual solver
inversion, the experience weight~\eqref{eq:weight}, and the team-from-members
composition~\eqref{eq:cfprior} follow the open-source
\texttt{calculate\_problem\_rating} logic in
\href{https://github.com/aropan/clist}{\texttt{github.com/aropan/clist}}. The
joint estimation in Chapter~5 is the Rasch / two-parameter item-response model
long used in educational testing, and the clock model in Chapter~6 is a
proportional-hazards survival model with right-censoring. The cross-region linking
argument is original to this project's setting but rests on the same
identifiability fact that underlies every paired-comparison rating system.

\end{document}